
\documentclass{article}
\usepackage{colortbl}
\usepackage{makecell}
\usepackage{multirow}
\usepackage{supertabular}

\begin{document}

\newcounter{utterance}

\centering \large Interaction Transcript for game `wordle', experiment `medium\_frequency\_words\_no\_clue\_no\_critic', episode 5 with reflexion{-}Qwen3.5{-}27B.
\vspace{24pt}

{ \footnotesize  \setcounter{utterance}{1}
\setlength{\tabcolsep}{0pt}
\begin{supertabular}{c@{$\;$}|p{.15\linewidth}@{}p{.15\linewidth}p{.15\linewidth}p{.15\linewidth}p{.15\linewidth}p{.15\linewidth}}
    \# & \multicolumn{2}{c}{Player} && \multicolumn{2}{c}{Game Master} \\
    \hline

    \theutterance \stepcounter{utterance}  
    & & & \multicolumn{4}{p{0.6\linewidth}}{
        \cellcolor[rgb]{0.9,0.9,0.9}{
            \makecell[{{p{\linewidth}}}]{
                \texttt{\tiny{[P1$\langle$GM]}}
                \texttt{You are a language wizard who likes to guess words by using the given rules.} \\
\\ 
\texttt{Welcome to Wordle! You have six attempts to guess the target word, a valid English word of five lowercase letters (a{-}z). Please use the tags "explanation:" and "guess:" to provide a concise explanation for each guess.} \\
\\ 
\texttt{For instance, if your guess is "apple", your response should be} \\
\texttt{explanation: this is a common five{-}letter English word, and I am starting my guess with this word.} \\
\texttt{guess: apple} \\
\\ 
\texttt{After each guess, your answer will be validated, and you will receive feedback indicating which letters are correct (green), which letters are correct but in the wrong position (yellow), and which letters are incorrect (red). This feedback can be useful in determining which letters to include or exclude in your next guess.} \\
\\ 
\texttt{For example, the feedback for "apple" might be:} \\
\texttt{guess\_feedback: a<yellow> p<yellow> p<green> l<yellow> e<red>} \\
\\ 
\texttt{The explanation should contain details about how the guess\_feedback is used to arrive at a new guess.} \\
\\ 
\texttt{Do not generate any other text.} \\
\\ 
\texttt{Let's begin with your first guess.} \\
            }
        }
    }
     \\ \\

    \theutterance \stepcounter{utterance}  
    & \multicolumn{4}{p{0.6\linewidth}}{
        \cellcolor[rgb]{0.561,0.737,0.561}{
            \makecell[{{p{\linewidth}}}]{
                \texttt{\tiny{[P1$\rangle$GM]}}
                \texttt{explanation: I am starting with a common five{-}letter word containing frequent vowels and consonants to maximize the information gained from the first feedback.} \\
\texttt{guess: slate} \\
            }
        }
    }
    & & \\ \\

    \theutterance \stepcounter{utterance}  
    & & & \multicolumn{2}{p{0.3\linewidth}}{
        \cellcolor[rgb]{0.9,0.9,0.9}{
            \makecell[{{p{\linewidth}}}]{
                \texttt{\tiny{[GM$|$GM]}}
                \texttt{attempts = 1} \\
\texttt{target\_word = steam} \\
\texttt{guess = slate} \\
\texttt{guess\_feedback = s<green> l<red> a<yellow> t<yellow> e<yellow>} \\
            }
        }
    }
    & & \\ \\

    \theutterance \stepcounter{utterance}  
    & & & \multicolumn{4}{p{0.6\linewidth}}{
        \cellcolor[rgb]{0.9,0.9,0.9}{
            \makecell[{{p{\linewidth}}}]{
                \texttt{\tiny{[P1$\langle$GM]}}
                \texttt{guess\_feedback: s<green> l<red> a<yellow> t<yellow> e<yellow>} \\
\\ 
\texttt{Provide your response only in this format.} \\
\texttt{explanation: details} \\
\texttt{guess: word} \\
            }
        }
    }
     \\ \\

    \theutterance \stepcounter{utterance}  
    & \multicolumn{4}{p{0.6\linewidth}}{
        \cellcolor[rgb]{0.561,0.737,0.561}{
            \makecell[{{p{\linewidth}}}]{
                \texttt{\tiny{[P1$\rangle$GM]}}
                \texttt{explanation: The feedback shows 's' is correct in the first position (green). 'l' is incorrect (red). 'a', 't', and 'e' are correct letters but in the wrong positions (yellow). I need to place 'a', 't', and 'e' in positions 2, 3, 4, or 5 while keeping 's' at the start and avoiding 'l'. The word "state" fits these constraints perfectly, placing 't' in position 2, 'a' in position 3, 't' again (which is allowed if the target has two t's, or testing the position), and 'e' at the end. Wait, "state" has two t's. Let's re{-}evaluate. Target has s, a, t, e. 'l' is out. Possible arrangement: s \_ \_ \_ \_. 'a', 't', 'e' need to move. If I guess "stale", 'l' is red so that's bad. How about "steal"? 'l' is red. "stare"? 'r' is new. "stave"? 'v' is new. "stale" is out. "state" uses s, t, a, t, e. If the target is "state", the feedback for "slate" would be s(green), l(red), a(yellow), t(yellow), e(yellow). This matches. But I need to be sure. Let's try "stare" to test 'r' and move 'a', 't', 'e'. Or "steal" is out. "stale" is out. "stave"? No. Let's try "stare". It keeps s, moves a, t, e, and adds r. Actually, "state" is a very strong candidate if the target has two t's. But "slate" only had one t. If the target is "state", the 't' in "slate" (pos 4) is yellow, meaning there is a 't' elsewhere. In "state", 't' is at 2 and 4. So 't' at 4 in "slate" would be yellow if the target is "state" (since the 't' at 2 is the match? No, usually the first match is green). If target is "state", and guess is "slate": s(green), l(red), a(yellow {-} target has a at 3), t(yellow {-} target has t at 2 and 4, guess has t at 4. If the algorithm matches the first t, it might be green? No, position 4 in target is t. So if guess is "slate", pos 4 is t. Target "state" pos 4 is t. That would be} \\
            }
        }
    }
    & & \\ \\

    \theutterance \stepcounter{utterance}  
    & & & \multicolumn{2}{p{0.3\linewidth}}{
        \cellcolor[rgb]{0.9,0.9,0.9}{
            \makecell[{{p{\linewidth}}}]{
                \texttt{\tiny{[GM$|$GM]}}
                \texttt{Error: The response should contain the 'guess:' keyword exactly once.} \\
            }
        }
    }
    & & \\ \\

    \theutterance \stepcounter{utterance}  
    & & & \multicolumn{2}{p{0.3\linewidth}}{
        \cellcolor[rgb]{0.9,0.9,0.9}{
            \makecell[{{p{\linewidth}}}]{
                \texttt{\tiny{[GM$|$GM]}}
                \texttt{game\_result = ABORT} \\
            }
        }
    }
    & & \\ \\

    \theutterance \stepcounter{utterance}  
    & & \multicolumn{4}{p{0.6\linewidth}}{
        \cellcolor[rgb]{0.95,0.95,0.95}{
            \makecell[{{p{\linewidth}}}]{
                \texttt{\tiny{[P1$\leftrightarrow$P]}}
                \texttt{Post{-}episode Reflexion (not sent to GM):} \\
\texttt{New plan: The previous attempt failed because I violated the rule "Do not generate any other text" by including a long, rambling internal monologue and incomplete reasoning within the "explanation" field, rather than providing a concise summary. Additionally, the response was cut off mid{-}sentence. Next time, I will strictly adhere to the format by providing only a brief, complete explanation of the logic (e.g., "s is fixed at pos 1; l is excluded; a, t, e are yellow and must be rearranged") followed immediately by a valid five{-}letter guess that respects the feedback constraints, ensuring no extra text or unfinished thoughts are included.} \\
            }
        }
    }
    & \\ \\

\end{supertabular}
}

\end{document}
